% Heap naming/style: daily-problems/10-3/CF1144B.cpp; two heaps here split by order.
% Practice: https://www.luogu.com.cn/problem/P1168
\subsection{对顶堆：动态维护中位数}
\label{ds:median-heap}
\AlgorithmInfo
    {连续插入数字，随时查询当前中位数及中间两个数。}
    {单次插入 $O(\log(n+1))$，查询 $O(1)$；空间 $O(n)$。}
    {保存 long long 值；查询前须非空，偶数个元素默认取较小的中位数。}
\subsubsection*{原理}
用大根堆 $q_1$ 保存较小的一半，用小根堆 $q_2$ 保存较大的一半，始终维护两个条件：$q_1$ 中任意元素不大于 $q_2$ 中任意元素，且
\[
 |q_1|=|q_2|\quad\text{或}\quad |q_1|=|q_2|+1.
\]
插入时与 $q_1$ 的堆顶比较，放入合适一侧；若大小失衡，将较多一侧的堆顶移到另一侧。每次插入后至多移动一个元素，顺序与大小条件都能保持。

因此，$q_1$ 恰好保存排序后的前 $\lceil n/2\rceil$ 个元素，它的堆顶就是较小的中位数。元素个数为奇数时只有一个中间数；为偶数时，两个堆顶分别是中间两个数。堆的比较器写法见第~\ref{stl:heap}~节。
\subsubsection*{使用教程}
\InlineCode{init()} 清空两堆；\InlineCode{add(x)} 插入一个值，\InlineCode{query()} 返回较小的中位数。\InlineCode{middle()} 返回中间两个数，奇数个元素时两项相同；\InlineCode{size()/empty()} 查询当前元素数和是否为空。新声明的对象本身为空，也可直接插入。
\begin{lstlisting}[style=mtycpp,numbers=none]
MedianHeap heap;
heap.init();
for(long long x:vector<long long>{5,1,3,3}) heap.add(x);
long long median=heap.query();
auto [l,r]=heap.middle();
int cnt=heap.size();
heap.init();
heap.add(-2);
\end{lstlisting}
\href{https://www.luogu.com.cn/problem/P1168}{P1168：中位数}在每次插入后，遇到奇数长度前缀就输出 \InlineCode{query()}。若题目规定偶数时取两个中间数的平均值，使用 \InlineCode{middle()}，相加前先转成 128 位或浮点类型。
\needspace{7\baselineskip}
\subsubsection*{模板代码}
\CodeFile{code/data-structures/median_heap.cpp}
\subsubsection*{需要注意}
本模板维护插入与中位数查询，不提供任意元素删除。维护固定第 $k$ 小时，可把左堆目标大小改成 $k$，仍保持左右有序，左堆顶即为答案；同时要保证当前元素数至少为 $k$，不能继续沿用本模板的两堆近似等大条件。
